\begin{abstract}
临床长文本包含丰富的诊疗过程信息，是表型识别、风险预测与临床辅助决策的重要信息来源。然而，这类文本通常具有篇幅长、术语密集、书写噪声强和标签边界模糊等特点，使标准确定性输出层难以同时兼顾判别性能与可信性表达。针对这一问题，本文提出一种面向临床长文本预测的可信输出层框架 \method{}。该方法以 \encoder{} 作为冻结式长上下文编码器，在特征空间中联合建模预测均值、输入相关方差以及基于方差诱导的加权最小二乘精修过程，从而将深度语义表示与结构化统计求解结合起来。本文围绕临床表型多标签识别任务，对 \method{} 的建模动机、形式化定义、两阶段训练策略和实现细节进行了系统阐述，并构建了线性头、MLP 头、Ridge 头与可信输出层之间的比较框架。阶段性实验结果表明，完整 \method{} 在 Phenotype 任务上可达到 0.786 的 weighted F1、0.804 的 micro F1 和 0.889 的 macro AUROC，同时在 MSE 和 MAE 上优于线性头、MLP 头以及仅采用阶段 A 训练的模型版本，显示出较好的综合性能与稳健性。本文的主要结论是：当临床长文本已经能够通过强编码器获得高质量表示时，围绕输出层开展可信建模仍然具有独立而明确的研究意义。
\end{abstract}
